\section*{Erd\H{o}s Problem 648: decreasing greatest prime factors}

\paragraph{Statement and claim.}
Write \(P(a)\) for the greatest prime factor of \(a\geq2\).
Let \(g(n)\) be the largest length of a sequence
\[
 2\leq a_1<\cdots<a_t<n,\qquad P(a_1)>\cdots>P(a_t).
\]
We reconstruct Cambie's determination of the order of magnitude:
\[
 (2-o(1))\sqrt{\frac n{\log n}}
 \ \leq g(n)\ \leq\
 (2\sqrt2+o(1))\sqrt{\frac n{\log n}}.
\]
The source is Stijn Cambie, \emph{On Erd\H{o}s problem \#648},
\url{https://arxiv.org/abs/2503.22691}; see also
\url{https://www.erdosproblems.com/648}.
The only asymptotic input below is the prime number theorem.
The existence or value of a limiting leading constant is not asserted.

\paragraph{Upper bound.}
Put \(p_i=P(a_i)\) and \(b_i=a_i/p_i\). Since \(a_{i+1}>a_i\) and
\(p_{i+1}<p_i\), we have \(b_{i+1}>b_i\). These are positive integers,
so \(b_i\geq i\). If \(q_j\) is the \(j\)-th prime, then the \(t-i+1\)
distinct primes \(p_i,\ldots,p_t\) show that \(p_i\geq q_{t-i+1}\).
Therefore
\[
 n>a_i=b_ip_i\geq i q_{t-i+1}.
\]
Taking \(i=\lfloor t/2\rfloor\), and using \(q_j\sim j\log j\), gives
\[
 n\geq(1-o(1))\frac{t^2\log t}{4}.
\]
The lower bound proved next shows that the extremal \(t=g(n)\)
is at least a constant times \(\sqrt{n/\log n}\).
The displayed upper bound then implies
\(\log t=(\tfrac12+o(1))\log n\), and hence the claimed
upper constant \(2\sqrt2\).

\paragraph{Lower bound.}
Fix \(0<\varepsilon<1\), and put
\(z=(1-\varepsilon)\sqrt{n\log n}\).
List all primes in \((\sqrt n,z]\) in decreasing order as
\(p_1>\cdots>p_r\). Starting with \(a_0=0\), define
\[
 a_i=p_i\left(\left\lfloor\frac{a_{i-1}}{p_i}\right\rfloor+1\right).
\]
Then \(a_i>a_{i-1}\) and
\[
 a_i\leq\sum_{j=1}^i p_j\leq\sum_{p\leq z}p
       \sim\frac{z^2}{2\log z}
       \sim(1-\varepsilon)^2n<n.
\]
Here the prime-sum estimate follows from the prime number theorem
by partial summation. Since \(p_i>\sqrt n\) and \(a_i<n\), the
cofactor \(a_i/p_i\) is smaller than \(\sqrt n<p_i\).
All its prime factors are therefore smaller than \(p_i\), proving
\(P(a_i)=p_i\). Finally,
\[
 r=\pi(z)-\pi(\sqrt n)
   =(2(1-\varepsilon)+o(1))\sqrt{\frac n{\log n}}.
\]
Letting \(\varepsilon\) tend to zero gives the lower bound.

\paragraph{Scope.}
This proves the known order-of-magnitude resolution, including the two
displayed asymptotic constants, rather than an exact formula for \(g(n)\).
It is an attributed reconstruction, not a new resolution.

\paragraph{Completion estimate.}
\textbf{COMPLETION: 100\%}
